% Brief (tikz-diagram-craft), register row ch4-20.
% Reader question: does a completion claim actually resolve against a
%   checkable predicate, or does it silently pass on nothing?
% Claim: a completion claim is CLAIMED, then a verifier runs the acceptance
%   predicate; a met predicate closes DONE, an unmet one is REFUSED, and an
%   unverifiable one must FAIL-LOUD rather than silently pass -- three
%   distinct outcomes, never a default accept.
% Evidence: Def. completionist, whitepaper/legible-swarm.tex, sec 4.4.5
%   (this pass); the swk commitment-oracle figure draws the same claim to
%   gate to DONE/reject shape this fragment reuses.
% Must distinguish: FAIL-LOUD (honest abstention, the mechanisms own point)
%   from REFUSED (a checked, failing predicate) -- collapsing them would
%   hide that fail-loud is correct behavior, not an error state.
% Grammar chosen: small state machine, one spine into a three-way guarded
%   gate, each branch a numbered guard keyed to a ruled legend.
% Grammar rejected: a two-outcome DONE/reject machine (the swk oracle
%   original) would silently merge unverifiable into refused, exactly the
%   silent-pass failure Def. completionist forbids naming.
% Five-second test: which of the three gates fires when the verifier itself
%   cannot run, and is that gate distinct from an outright failing check?
\begin{figure}[H]
\centering
\pdfigurehue{pdteal}%
\begin{tikzpicture}[pd figure,x=1cm,y=1cm]
  \def\lx{1.30}   % claim column
  \def\rx{8.60}   % outcome column
  \node[pd focus state,minimum width=2.5cm] (C) at (\lx,1.60) {claimed};
  \node[pd ready state,minimum width=3.0cm]  (D) at (\rx,3.10) {done};
  \node[pd warn state,minimum width=3.0cm]   (U) at (\rx,1.60) {fail-loud};
  \node[pd breach state,minimum width=3.0cm] (R) at (\rx,0.10) {refused};
  \draw[pd focus arrow] (C.east) -- (4.95,1.60) -- (D.west);
  \draw[pd focus arrow] (C.east) -- (U.west);
  \draw[pd focus arrow] (C.east) -- (4.95,1.60) -- (R.west);
  \node[pd title,anchor=south] at (3.10,1.90) {run verifier};
  \node[pd badge] at (6.55,2.62) {1};
  \node[pd badge] at (6.55,1.60) {2};
  \node[pd badge] at (6.55,0.58) {3};
  \def\ly{-1.55}
  \draw[pd rule] (0,\ly) -- (10.9,\ly);
  \node[pd title,anchor=north west] at (0.05,{\ly-.35}) {guard};
  \node[pd title,anchor=north west] at (5.30,{\ly-.35}) {what breaks without it};
  \draw[pd hairline] (0,{\ly-1.00}) -- (10.9,{\ly-1.00});
  \foreach \i/\n/\g/\c in {0/1/{predicate met}/{a hollow claim reads as done},
    1/2/{predicate unverifiable}/{an unverifiable claim silently auto-passes},
    2/3/{predicate unmet}/{a failing claim is accepted anyway}}{
    \node[pd badge] at (.28,{\ly-1.35-\i*.60}) {\n};
    \node[pd label,anchor=west,align=left] at (.62,{\ly-1.35-\i*.60}) {\g};
    \node[pd label,anchor=west,align=left,text=pderror] at (5.30,{\ly-1.35-\i*.60}) {\c};}
\end{tikzpicture}
\caption{The completionist obligation of Def.~\ref{def:completionist}: a
completion claim is CLAIMED, then the authority runs a verifier derived
from the tasks own acceptance criteria. A met predicate closes done
(green); an unmet predicate is refused (red) rather than accepted anyway;
and a predicate the verifier cannot evaluate must fail loudly (amber) and
surface I cannot verify this completion rather than silently passing.
Fail-loud is a third, distinct outcome from refused: it is the mechanism
working as designed, not a failure the way refused is. [internal]}
\label{fig:completion-verifier}
\end{figure}
